\section{Judge Prompt}
\label{app:prompt}

Both LLM judges receive the identical step-grading rubric and template below
(from \texttt{src/llm\_judge.py}); only the backend model differs. The judge is
called at temperature $0$ and sees the question, the rendered prior steps, and
the current step's expression and operation. Verdicts are parsed from a strict
JSON object; unparseable responses are treated as \textsc{abstain}.

\paragraph{Step rubric.}
\begin{quote}\small\ttfamily
You are a strict grader for a single step of a math derivation. You will be
given the question being solved, the preceding steps (as context), and the
current step's symbolic expression with the operation connecting it to the
previous step.\\[4pt]
Decide whether the current step is mathematically valid given the previous
step --- that is, does the stated operation, applied to the previous expression,
produce this expression?\\[4pt]
Rules:\\
- "passed": the step is a valid algebraic, calculus, or definitional move.\\
- "failed": the step contains a math error (wrong manipulation, wrong limits,
wrong sign, dropped term, etc.).\\
- "abstain": the step is too vague, uses undefined symbols, or you cannot tell
without more context. Prefer "abstain" over guessing.\\[4pt]
Output ONLY a JSON object in this exact shape, no prose before or after:\\
\{ "verdict": "passed" | "failed" | "abstain", "confidence": <float in [0,1]>,
"reasoning": "<one or two sentence justification>" \}
\end{quote}

\paragraph{Step template.}
\begin{quote}\small\ttfamily
\{rubric\}\\[4pt]
QUESTION: \{question\}\\[4pt]
PRIOR STEPS (most recent last): \{prior\_steps\}\\[4pt]
CURRENT STEP:\\
\hspace*{1em}expression: \{expression\}\\
\hspace*{1em}operation connecting to previous: \{operation\}
\end{quote}

\section{Injection Mutation Catalog}
\label{app:mutations}

Table~\ref{tab:mutations} lists the 14 mutation operators across five families
(from \texttt{src/error\_injection.py}). Three families inject genuine errors
(severity \emph{error}); the \emph{equivalent} and \emph{control} families are
negatives that a correct grader must not flag. First steps (definitions) are not
mutated. Mutations that do not apply to a given expression produce a
\emph{skipped} trial that is excluded from all metrics.

\begin{table*}[tp]
\centering
\small
\begin{tabular}{lll}
\toprule
family & mutation & effect \\
\midrule
\multirow{7}{*}{algebraic}
 & sign\_flip & negate leading sign \\
 & coeff\_x2 & multiply by 2 \\
 & coeff\_half & divide by 2 \\
 & drop\_term & drop last additive term \\
 & exponent\_bump & increment first integer exponent \\
 & constant\_shift & add 1 \\
 & missing\_factor\_of\_2 & drop leading factor of 2 \\
\midrule
\multirow{2}{*}{bounds}
 & swap\_integral\_limits & swap $a\leftrightarrow b$ in first integral/sum \\
 & drop\_upper\_limit & corrupt upper limit \\
\midrule
identifier & rename\_variable & rename a key symbol in this step \\
\midrule
\multirow{3}{*}{equivalent}
 & commute\_operands & swap operands of first $+$/$*$ \\
 & expand\_power & rewrite $X^2$ as $X\cdot X$ \\
 & distribute & expand $k(a{+}b)$ into $ka{+}kb$ \\
\midrule
control & no\_op & identity (no mutation) \\
\bottomrule
\end{tabular}
\caption{The 14 mutation operators. \emph{algebraic}, \emph{bounds}, and
\emph{identifier} are severity \emph{error}; \emph{equivalent} is severity
\emph{equivalent}; \emph{control} is severity \emph{control}. Graded trial
counts per family (after skips): $74$ algebraic, $6$ bounds, $14$ identifier,
$10$ equivalent, $14$ control ($118$ total).}
\label{tab:mutations}
\end{table*}

\section{Human Error Detectability}
\label{app:human-errors}

Table~\ref{tab:human-errors} classifies each of the 13 human-flagged errors as
\emph{in-fragment} (symbolically detectable in principle by an ideal CAS step
checker) or \emph{out-of-fragment} (definition, notation-strictness, or
incomplete presentation). Criterion: an error is in-fragment iff algebraic
equivalence against a correct continuation of the prior math could refute the
written expression. Out-of-fragment errors include bilateral-vs-unilateral
Laplace conventions, wrong signal definitions relative to the intended question,
odd/even characterizations tied to that setup, IBP boundary terms omitted when
both limits evaluate to zero, and incomplete equations under a notation-strict
rubric. Result: $9/13$ out-of-fragment; $4/13$ in-fragment, of which the
deployed verifier caught $0$ (2 false passes, 2 abstains).

Table~\ref{tab:human-errors} distinguishes two ways the verifier commits PASS
on an error step. A \emph{false pass} (5 steps) is an active wrong equivalence
judgment: the verifier tests the step against a prior expression and wrongly
finds it consistent. A \emph{pass (1st)} (2 steps, Q03.1 and Q10.1) is a
vacuous pass: the step is the first in its derivation, so there is no prior
expression to test equivalence against, and the verifier cannot fail by
construction. Both count as committed passes on an error step, so $5+2=7$
matches the symbolic row's false-negative count in Table~\ref{tab:threeway}.
\S\ref{sec:cascade} spells out this split at the point where the count is
used; the abstract and \S\ref{sec:intro} use ``false passes'' more loosely, as
shorthand for all 7 committed passes on error steps.

\begin{table*}[tp]
\centering
\footnotesize
\setlength{\tabcolsep}{3pt}
\begin{tabular}{lllll}
\toprule
step & type & fragment & verifier & reason (abbrev.) \\
\midrule
Q03.1 & definition & out & pass (1st) & bilateral vs unilateral LT \\
Q09.1 & definition & out & abstain & wrong $\pm 1$ square vs $0/1$ pulse \\
Q09.3 & algebra & \textbf{in} & false pass & wrong antiderivative limits \\
Q09.4 & algebra & \textbf{in} & false pass & wrong lower-limit exponential \\
Q09.6 & algebra & \textbf{in} & abstain & carries Q09.3--4; wrong $\pi$ vs $\pi/2$ \\
Q09.7 & missing\_term & \textbf{in} & abstain & drops second integral term \\
Q09.11 & definition & out & false pass & odd/even claim vs intended signal \\
Q10.1 & definition & out & pass (1st) & bilateral vs unilateral LT \\
Q16.2 & missing\_term & out & abstain & IBP $\infty$ term omitted; both $0$ \\
Q22.5 & missing\_term & out & abstain & LHS only; missing ${=}X(s)$ \\
Q22.7 & missing\_term & out & false pass & incomplete equation fragment \\
Q22.8 & missing\_term & out & abstain & ratio symbol; no RHS \\
Q22.9 & missing\_term & out & false pass & divide claimed; RHS still missing \\
\bottomrule
\end{tabular}
\caption{All 13 human-labeled errors. ``fragment'' $=$ in/out of the symbolic
equivalence fragment; ``verifier'' $=$ deployed SymPy outcome on that step.
In-fragment count: $4$; out-of-fragment: $9$; caught by verifier: $0/13$.}
\label{tab:human-errors}
\end{table*}

\section{Qualitative Examples}
\label{app:examples}

Table~\ref{tab:examples} shows four qualitative failure modes referenced in
\S\ref{sec:analysis}, each a concrete frozen step with truncated judge
rationale.

\begin{table*}[tp]
\centering
\footnotesize
\setlength{\tabcolsep}{3pt}
\begin{tabular}{p{1.55cm}p{3.6cm}ccc p{4.4cm}}
\toprule
mode & step (abbrev.) & sym & gpt & claude & note / rationale (trunc.) \\
\midrule
Equiv.\ rewrite (judge FP)
 & $2\sin\omega\cos\omega \!\to\! \sin\omega\cos\omega\cdot 2$
   (injection, commute)
 & pass & pass & \textbf{fail}
 & Claude: ``non-standard order\ldots\ potential misunderstanding of the
   identity''; algebraically equal. \\
\midrule
Semantic / definition (sym miss)
 & Q03.1: $\int_{-\infty}^{\infty}x(t)e^{-st}\,dt$
   (bilateral LT as start)
 & \textbf{pass} & pass & fail
 & Human: unilateral LT expected. Claude: ``defined from $0$ to $\infty$,
   not $-\infty$ to $\infty$.'' \\
\midrule
In-fragment algebra (sym false pass)
 & Q09.3: $\frac{1-e^{-Ik\omega_0 T_0/2}}{Ik\omega_0 T_0}$
   (wrong integral eval.)
 & \textbf{pass} & fail & fail
 & Limits $T_0/4$ written as $T_0/2$. Both judges refute; cascade would
   \emph{inherit} the symbolic pass. \\
\midrule
Bounds (sym abstain, judge catch)
 & $\int_0^{\infty}\!\to\!\int_{\infty}^{0}$
   (injection, swap limits)
 & abstain & \textbf{fail} & fail
 & GPT: ``limits reversed\ldots\ without a minus sign.'' Symbolic silent;
   judge covers the family. \\
\bottomrule
\end{tabular}
\caption{Four qualitative failure modes from the frozen artifacts. Bold =
the distinctive (wrong or uniquely useful) verdict for that row. Judge text is
verbatim but truncated; full rationales are in the released JSON.}
\label{tab:examples}
\end{table*}

\section{Cost and Latency Detail}
\label{app:cost}

LLM wall-clock RTT was not logged in the frozen runs, so we price the cascade
of \S\ref{sec:results} by \emph{call count} and estimated tokens.
Reconstructing the released judge prompt yields $\approx 419$ input /
$\approx 58$ output tokens per human step (output from returned JSON only;
hidden reasoning tokens for \texttt{gpt-oss} are excluded, so its dollar
figures are lower bounds). At published list rates (\texttt{gpt-oss-120b}:
$\$0.35/\$0.75$ per MTok in/out on Cerebras; \texttt{claude-sonnet-4-5}:
$\$3/\$15$ on Anthropic), grading all 191 human steps with the judge alone
costs about $\$0.036$ (\texttt{gpt-oss}) or $\$0.41$ (\texttt{claude}). The
cascade cuts LLM calls from 191 to 101 on the human set ($47\%$ fewer; about
$\$0.019$ / $\$0.22$), because symbolic already commits on 90 steps, and from
118 to 76 on the 118 injection trials ($36\%$ fewer). So the injection
precision win of \S\ref{sec:cascade} arrives with roughly a third fewer paid
calls, and every symbolic commit replaces an API round trip with local CPU.
Prices change; what stays fixed is the call-count reduction, which any
deployment can measure without our rate card.

\section{Reproduction}
\label{app:repro}

All artifacts are regenerated with fixed scripts; the symbolic grader is
deterministic and the judges are called at temperature $0$ with per-prompt
caching. From \texttt{<repo-root>/} (with \texttt{PYTHONPATH=.}):

\begin{quote}\small\ttfamily
\# judge vs human (all 191 steps)\\
python scripts/judge\_vs\_human.py\\[2pt]
\# paired McNemar + Holm--Bonferroni summary\\
python scripts/build\_paired\_stats\_summary.py\\[2pt]
\# error injection with judge\\
python -m src.error\_injection \textbackslash\\
\hspace*{1em}-{}-with-judge\\[2pt]
\# tables + figures\\
python scripts/build\_threeway\_table.py\\
python scripts/build\_cascade\_table.py\\
python scripts/build\_confidence\_curves.py\\
python scripts/build\_judge\_agreement.py\\
python scripts/build\_cost\_latency.py\\
python scripts/render\_figures.py\\[2pt]
\# second judge (Claude):\\
\hspace*{1em}JUDGE\_BACKEND=anthropic + separate output files
\end{quote}

The evaluator emits a schema-v2 reproducibility envelope (model, seeds, config)
alongside each result file so runs can be matched to their exact configuration.
This appendix lists only scripts used to produce the reported tables and figures;
we do not report inter-annotator agreement, and no dual-annotation pipeline is
part of the reproduction path for this submission. An anonymous archive with
code, prompts, labels, and frozen artifacts is at \url{https://osf.io/5kvqe/?view_only=8c0032bd76244a1daaebf6cb3feb2c92}.
